% Part: history
% Chapter: biographies
% Section: bertrand-russell

\documentclass[../../../include/open-logic-section]{subfiles}

\begin{document}

\olfileid{his}{bio}{rus} 

\olsection{బెర్‌ట్రాండ్ రసెల్}

\olphoto{russell-bertrand}{బెర్‌ట్రాండ్ రసెల్}

ఆధునిక విశ్లేషణాత్మక తత్వశాస్త్ర స్థాపకులలో ఒకరిగా బెర్‌ట్రాండ్ రసెల్ ప్రశంసలు పొందారు. 1872 మే 18న జన్మించిన రసెల్ తత్వశాస్త్రం, తర్కశాస్త్రంలో చేసిన కృషికే కాక వివిధ రంగాలలో ప్రజాదరణ పొందిన ఎన్నో పుస్తకాలు రాసినందుకూ ప్రసిద్ధి. జీవితాంతం రాజకీయ ఉద్యమాల్లోనూ ఉత్సాహంగా పాల్గొన్నారు.

రసెల్ వేల్స్‌లోని మాన్‌మత్‌షైర్‌లో ట్రెలెక్‌లో జన్మించారు. ఆయన తల్లిదండ్రులు బ్రిటిష్ కులీన వర్గానికి చెందినవారు. వారు స్వతంత్ర ఆలోచనాపరులు; అప్పటి బోస్టన్‌లోని తీవ్ర సంస్కరణవాదులతో కూడా స్నేహం చేశారు. దురదృష్టవశాత్తూ, రసెల్ చిన్నతనంలోనే తల్లిదండ్రులు మరణించారు; ఆయన తాతయ్య, నానమ్మల వద్దకు వెళ్లాల్సి వచ్చింది. అక్కడ ఆయనను మతపద్ధతిలో పెంచారు (దాన్ని ఆయన తల్లిదండ్రులు ఎట్టి పరిస్థితుల్లోనూ కోరుకోలేదు). నైతిక విషయాలన్నిటిలోనూ ఆయన నానమ్మ చాలా కఠినంగా ఉండేవారు. కౌమార దశలో ఎక్కువగా వ్యక్తిగత బోధకుల ద్వారా ఇంట్లోనే చదువుకున్నారు.

విశ్లేషణాత్మక తత్వశాస్త్రం, ముఖ్యంగా తర్కశాస్త్రంపై రసెల్ ప్రభావం అపారం. కేంబ్రిడ్జ్‌లోని ట్రినిటీ కళాశాలలో గణితం, తత్వశాస్త్రం చదివారు; అక్కడ గణితశాస్త్రజ్ఞుడు, తత్వవేత్త ఆల్ఫ్రెడ్ నార్త్ వైట్‌హెడ్ ప్రభావం ఆయనపై పడింది. 1910లో రసెల్, వైట్‌హెడ్ \emph{Principia Mathematica} తొలి సంపుటాన్ని ప్రచురించారు; అందులో గణితాన్ని తర్కశాస్త్రానికి తగ్గించవచ్చనే అభిప్రాయాన్ని సమర్థించారు. తరువాత రసెల్ వందలాది పుస్తకాలు, వ్యాసాలు, రాజకీయ కరపత్రాలు ప్రచురించారు. 1950లో సాహిత్య నోబెల్ బహుమతి పొందారు.

రాజకీయాలు, సామాజిక ఉద్యమాల్లో రసెల్ లోతుగా నిమగ్నమయ్యారు. మొదటి ప్రపంచ యుద్ధం సమయంలో యుద్ధ వ్యతిరేక కార్యకలాపాలు, నిరసనల కారణంగా ఆయనను అరెస్టు చేసి ఆరు నెలలు జైలుకు పంపారు. జైలులో ఉండగా చదవగలిగారు, రాయగలిగారు; ఆ అనుభవం ``చాలా ఆహ్లాదకరంగా'' అనిపించిందని చెప్పారు. జీవితాంతం శాంతివాదిగా ఉన్నారు; 1961లో అణు నిరాయుధీకరణ సభకు హాజరైనందుకు మళ్లీ కారాగారంలో ఉంచారు. 1948లో జరిగిన విమాన ప్రమాదం నుంచి కూడా బయటపడ్డారు; ఆ ప్రమాదంలో ధూమపాన విభాగంలో కూర్చున్నవారే బతికారు. అందుచేత ధూమపానమే తన ప్రాణాలను కాపాడిందని రసెల్ చెప్పుకున్నారు. నాలుగుసార్లు వివాహం చేసుకున్న రసెల్‌కు వివాహేతర సంబంధాలు కొనసాగించేవారని పేరుంది. 1970 ఫిబ్రవరి 2న వేల్స్‌లోని పెన్‌రిన్‌డెయుడ్రెత్‌లో 97 ఏళ్ల వయసులో మరణించారు.

\begin{reading}
రసెల్ 1872--1967 మధ్య తన జీవితాన్ని వివరించే మూడు భాగాల ఆత్మకథ రాశారు \citep{Russell1967,Russell1968,Russell1969}. మెక్‌మాస్టర్ విశ్వవిద్యాలయంలోని బెర్‌ట్రాండ్ రసెల్ పరిశోధనా కేంద్రంలో రసెల్ ఆర్కైవ్ ఉంది. ఆయన సంకలిత రచనల సంపుటాలు (శోధించదగిన సూచికలతో సహా), ఆర్కైవ్ ప్రాజెక్టుల గురించి ఆ కేంద్రం వెబ్‌సైట్‌కు \citet{Duncan2015} చూడండి. రసెల్ రాసిన \emph{On
  Denoting} \citep{Russell1905} ఇరవయ్యో శతాబ్దపు విశ్లేషణాత్మక తత్వశాస్త్రంలో క్లాసిక్ వ్యాసం.

స్టాన్‌ఫర్డ్ ఎన్‌సైక్లోపీడియా ఆఫ్ ఫిలాసఫీలోని రసెల్ వ్యాసంలో \citep{Irvine2015}, కోరిక, రాజకీయ సిద్ధాంతం గురించి రసెల్ మాట్లాడిన ధ్వని భాగాలు ఉన్నాయి. ఆయనతో జరిగిన అనేక వీడియో ముఖాముఖిలు ఆన్‌లైన్‌లో లభిస్తాయి. ఉదాహరణకు ధూమపానం, విమాన ప్రమాదం గురించి ఆయన మాట్లాడేది చూడాలంటే \citet{RussellND} చూడండి. \emph{Introduction to Mathematical Philosophy} సహా ఆయన కొన్ని రచనలు \citet{LibriVoxND}లో ఉచిత ఆడియో పుస్తకాలుగా అందుబాటులో ఉన్నాయి.
\end{reading}

\end{document}
